\section{Introducción}
Las herramientas digitales hoy en día forman parte de la vida cotidiana para múltiples actividades académicas, como:

\begin{multicols}{2}
    \begin{itemize}
        \item  Consultar bibliografía.
        \item Acceder a aulas virtuales.
        \item Elaborar documentos.
        \item Comunicarse con docentes y compañeros.
        \item Organizar tareas.
        \item Producir contenidos.
    \end{itemize}
\end{multicols}
 
Sin embargo, la disponibilidad de una herramienta no garantiza la finalidad académica clara para que todos los estudiantes tengan las mismas condiciones equivalentes de acceso, habilidades o acompañamiento. Así que, resulta pertinente describir cómo se utilizan estas herramientas en un contexto universitario concreto.

Esta propuesta delimita el estudio a estudiantes universitarios vinculados a instituciones de educación
superior ubicadas en Bogotá, Colombia. Se propone identificar qué herramientas digitales utilizan, con qué
frecuencia, para cuáles actividades académicas y cómo valoran su utilidad. El proyecto no presupone que la
tecnología mejore automáticamente el aprendizaje; busca producir una descripción ordenada de las prácticas
reportadas por los participantes.

El contexto nacional ayuda a explicar la relevancia del tema, aunque no sustituye la información específica de la población universitaria. El Departamento Administrativo Nacional de Estadística \parencite{dane2025entic} reportó que en 2024 el 82,7\% de los hogares de Bogotá tenía conexión a internet, frente al 65,6\% del total nacional. Estas cifras describen hogares y conectividad, no estudiantes universitarios ni uso educativo; por tanto, se emplean únicamente como indicador contextual de acceso.

La \textcite{unesco2023tecnologia} advierte que el potencial educativo de la tecnología depende de cómo se integra a la enseñanza y de las condiciones en que se utiliza. En consecuencia, estudiar la frecuencia, las finalidades y las percepciones de uso puede ofrecer información útil para orientar actividades de alfabetización digital y estrategias de apoyo académico, sin confundir asociación, percepción y causalidad.

\subsection{Delimitación}
La propuesta se centra en estudiantes mayores de edad, matriculados en programas técnicos profesionales, tecnológicos o universitarios de instituciones de educación superior de Bogotá durante el periodo de aplicación que se defina. Como todavía no se ha seleccionado una institución ni se dispone de un marco muestral, el alcance territorial es Bogotá y la selección definitiva de participantes queda condicionada a los permisos y recursos disponibles.

\subsection{Propósito del documento}
Presentar un diseño metodológico viable para una investigación descriptiva, establecer variables
observables y proponer procedimientos de recolección y análisis que puedan ejecutarse después de la aprobación
académica y ética correspondiente.